\documentclass[12pt, a4paper, twoside]{article}

% ── Font e codifica ──────────────────────────────────────────────────────────
\usepackage[T1]{fontenc}
\usepackage[utf8]{inputenc}
\usepackage{lmodern}
\usepackage{microtype}

% ── Lingua ───────────────────────────────────────────────────────────────────
\usepackage[italian]{babel}

% ── Matematica ───────────────────────────────────────────────────────────────
\usepackage{amsmath, amssymb, amsthm, mathtools}
\usepackage{physics}       % \dd, \pdv, \grad, \curl, ecc.
\usepackage{tensor}        % notazione tensoriale
\usepackage{bm}            % grassetto matematico

% ── Layout ───────────────────────────────────────────────────────────────────
\usepackage[
  a4paper,
  top=2.8cm, bottom=2.8cm,
  left=3.0cm, right=2.5cm,
  headheight=15pt
]{geometry}
\usepackage{fancyhdr}
\usepackage{emptypage}
\usepackage{titlesec}
\usepackage{setspace}
\onehalfspacing

% ── Grafica e colori ─────────────────────────────────────────────────────────
\usepackage{xcolor}
\usepackage{graphicx}
\usepackage{tikz}
\usepackage{pgfplots}
\pgfplotsset{compat=1.18}
\usetikzlibrary{arrows.meta, decorations.pathmorphing, patterns, calc, 3d}

% ── Codice sorgente ──────────────────────────────────────────────────────────
\usepackage{listings}
\usepackage{minted}         % syntax highlighting con Pygments

% ── Box e ambienti ───────────────────────────────────────────────────────────
\usepackage{tcolorbox}
\tcbuselibrary{theorems, skins, breakable, listings}
\usepackage{mdframed}

% ── Tabelle ──────────────────────────────────────────────────────────────────
\usepackage{booktabs}
\usepackage{multirow}
\usepackage{array}

% ── Hyperlink ────────────────────────────────────────────────────────────────
\usepackage[
  colorlinks=true,
  linkcolor=NavyBlue,
  citecolor=ForestGreen,
  urlcolor=RoyalBlue,
  pdftitle={Simulazione di un Buco Nero},
  pdfauthor={Trattazione Fisico-Matematica},
  bookmarks=true
]{hyperref}
\usepackage{cleveref}

% ── Bibligrafia ──────────────────────────────────────────────────────────────
\usepackage[numbers, sort&compress]{natbib}

% ─────────────────────────────────────────────────────────────────────────────
%  COLORI PERSONALIZZATI
% ─────────────────────────────────────────────────────────────────────────────
\definecolor{NavyBlue}{HTML}{1B3A6B}
\definecolor{RoyalBlue}{HTML}{2E5BCC}
\definecolor{ForestGreen}{HTML}{1A6B3A}
\definecolor{CrimsonRed}{HTML}{C0392B}
\definecolor{DarkGray}{HTML}{333333}
\definecolor{LightBlue}{HTML}{EAF0FB}
\definecolor{LightGold}{HTML}{FEF9E7}
\definecolor{CodeBg}{HTML}{F4F6F9}
\definecolor{BoxBorder}{HTML}{1B3A6B}

% ─────────────────────────────────────────────────────────────────────────────
%  AMBIENTI PERSONALIZZATI
% ─────────────────────────────────────────────────────────────────────────────

% Box per le equazioni fondamentali
\tcbset{
  equationbox/.style={
    enhanced, breakable,
    colback=LightBlue, colframe=NavyBlue,
    arc=3pt, boxrule=1.2pt,
    fonttitle=\bfseries\color{NavyBlue},
    attach boxed title to top left={yshift=-2mm, xshift=6mm},
    boxed title style={colback=NavyBlue, colframe=NavyBlue, arc=2pt}
  },
  curiositybox/.style={
    enhanced, breakable,
    colback=LightGold, colframe=CrimsonRed,
    arc=3pt, boxrule=1pt,
    fonttitle=\bfseries\color{CrimsonRed},
    attach boxed title to top left={yshift=-2mm, xshift=6mm},
    boxed title style={colback=CrimsonRed, colframe=CrimsonRed, arc=2pt}
  },
  codebox/.style={
    enhanced, breakable,
    colback=CodeBg, colframe=DarkGray,
    arc=2pt, boxrule=0.8pt,
    fonttitle=\bfseries\ttfamily\small\color{DarkGray},
  }
}

\newenvironment{physbox}[1]{%
  \begin{tcolorbox}[equationbox, title={#1}]
}{%
  \end{tcolorbox}
}

\newenvironment{curiosity}[1]{%
  \begin{tcolorbox}[curiositybox, title={Curiosit\`a: #1}]
}{%
  \end{tcolorbox}
}

% Teorema, definizione, osservazione
\theoremstyle{definition}
\newtheorem{defn}{Definizione}[section]
\newtheorem{prop}{Proposizione}[section]
\newtheorem{oss}{Osservazione}[section]

% ─────────────────────────────────────────────────────────────────────────────
%  INTESTAZIONI E PIÈ DI PAGINA
% ─────────────────────────────────────────────────────────────────────────────
\pagestyle{fancy}
\fancyhf{}
\fancyhead[LE]{\small\color{NavyBlue}\textit{Simulazione di un Buco Nero}}
\fancyhead[RO]{\small\color{NavyBlue}\textit{\leftmark}}
\fancyfoot[C]{\small\color{DarkGray}\thepage}
\renewcommand{\headrulewidth}{0.4pt}
\renewcommand{\headrule}{\hbox to\headwidth{\color{NavyBlue}\leaders\hrule height \headrulewidth\hfill}}

% ─────────────────────────────────────────────────────────────────────────────
%  STILE SEZIONI
% ─────────────────────────────────────────────────────────────────────────────
\titleformat{\section}
  {\color{NavyBlue}\Large\bfseries}
  {\color{NavyBlue}\thesection.}{0.8em}{}
  [\color{NavyBlue}\titlerule]

\titleformat{\subsection}
  {\color{NavyBlue}\large\bfseries}
  {\color{NavyBlue}\thesubsection.}{0.7em}{}

\titleformat{\subsubsection}
  {\color{CrimsonRed}\normalsize\bfseries}
  {\color{CrimsonRed}\thesubsubsection.}{0.6em}{}

% ─────────────────────────────────────────────────────────────────────────────
%  MACRO UTILI
% ─────────────────────────────────────────────────────────────────────────────
\newcommand{\RS}{r_{\!s}}          % raggio di Schwarzschild
\newcommand{\Rph}{r_{\!\text{ph}}}  % raggio fotosfera
\newcommand{\Risco}{r_{\!\text{ISCO}}} % ISCO
\newcommand{\bc}{b_{\!c}}          % parametro d'impatto critico
\newcommand{\Lobs}{L_{\text{obs}}} % luminosità osservata
\newcommand{\Lem}{L_{\text{em}}}   % luminosità emessa
\newcommand{\Tobs}{T_{\text{obs}}} % temperatura osservata
\newcommand{\Tem}{T_{\text{em}}}   % temperatura emessa
\newcommand{\munu}{{\mu\nu}}
\newcommand{\mink}{\eta_{\munu}}
\newcommand{\Gmunu}{G_{\munu}}
\newcommand{\Rmunu}{R_{\munu}}
\newcommand{\Tmunu}{T_{\munu}}
\newcommand{\gmunu}{g_{\munu}}

% ═════════════════════════════════════════════════════════════════════════════
%  DOCUMENTO
% ═════════════════════════════════════════════════════════════════════════════
\begin{document}

% ─────────────────────────────────────────────────────────────────────────────
%  FRONTESPIZIO
% ─────────────────────────────────────────────────────────────────────────────
\begin{titlepage}
  \centering
  \vspace*{2cm}

  {\color{NavyBlue}\rule{\linewidth}{1.5pt}}\\[0.5em]
  {\color{NavyBlue}\rule{\linewidth}{0.4pt}}\\[2em]

  {\Huge\bfseries\color{NavyBlue} Simulazione di un Buco Nero}\\[0.8em]
  {\large\itshape\color{DarkGray}
    Dalla Relativit\`a Generale al Ray Tracing su GPU}\\[2em]

  {\color{NavyBlue}\rule{\linewidth}{0.4pt}}\\[0.5em]
  {\color{NavyBlue}\rule{\linewidth}{1.5pt}}\\[3em]

  \begin{minipage}{0.85\textwidth}
    \centering
    {\normalsize\color{DarkGray}
      Trattazione fisico-matematica completa della metrica di Schwarzschild,
      delle geodetiche nulle, del lensing gravitazionale, del disco di
      accrescimento con profilo di Novikov--Thorne, del redshift gravitazionale
      e dell'effetto Doppler relativistico. Implementazione dello shader GLSL
      con integrazione numerica Runge--Kutta del quarto ordine.
    }
  \end{minipage}

  \vfill

  \begin{tabular}{ll}
    \toprule
    \textbf{Argomento}   & Relativit\`a Generale, Astrofisica, Computer Graphics \\
    \textbf{Livello}     & Laurea Magistrale in Fisica \\
    \textbf{Linguaggio}  & GLSL 3.30 (OpenGL Shading Language) \\
    \textbf{Unità}       & Geometrizzate: $G = c = M = 1$ \\
    \bottomrule
  \end{tabular}

  \vspace{2em}
  {\small\color{DarkGray}\today}
\end{titlepage}

% ─────────────────────────────────────────────────────────────────────────────
%  INDICE
% ─────────────────────────────────────────────────────────────────────────────
\tableofcontents
\newpage

% ═════════════════════════════════════════════════════════════════════════════
%  INTRODUZIONE
% ═════════════════════════════════════════════════════════════════════════════
\section*{Introduzione}
\addcontentsline{toc}{section}{Introduzione}

Un buco nero è forse l'oggetto più estremo che la fisica contempli: una regione
dello spaziotempo dalla cui influenza gravitazionale nulla può sfuggire, nemmeno
la luce. Eppure i buchi neri non sono completamente invisibili: il materiale che
cade verso di essi forma un \emph{disco di accrescimento} incandescente, e la
loro gravità deforma la geometria della luce che ci arriva, producendo immagini
straordinarie visibili anche da osservatori lontani.

Questo documento sviluppa la fisica e la matematica necessarie a \emph{simulare}
questa visione, con un livello di accuratezza sufficiente a riprodurre
qualitativamente le immagini ottenute da telescopi spaziali e da interferometri
radio come l'\textit{Event Horizon Telescope} (EHT).

Il percorso è il seguente:
\begin{enumerate}
  \item Le equazioni di campo di Einstein e la loro soluzione per un buco nero
        sferico (metrica di Schwarzschild).
  \item La soluzione per un buco nero in rotazione (metrica di Kerr) e
        le sue proprietà peculiari.
  \item Le geodetiche nulle: come si muove la luce in uno spaziotempo curvo.
  \item Il disco di accrescimento: struttura, temperatura, luminosità.
  \item Gli effetti osservativi: redshift gravitazionale, effetto Doppler
        relativistico, beaming.
  \item L'algoritmo di ray tracing relativistico e la sua implementazione
        in GLSL con Runge--Kutta al quarto ordine.
\end{enumerate}

\medskip
\noindent\textbf{Unità.} In tutto il documento utilizziamo \emph{unità
geometrizzate}: $G = c = 1$. In queste unità la massa $M$ ha le dimensioni
di una lunghezza (in unità SI, $G M / c^2$ è la metà del raggio di
Schwarzschild). Il raggio di Schwarzschild diventa semplicemente $\RS = 2M$.

\newpage

% ═════════════════════════════════════════════════════════════════════════════
%  PARTE I — FONDAMENTI FISICI
% ═════════════════════════════════════════════════════════════════════════════

\section{La Relativit\`a Generale e la Curvatura dello Spaziotempo}

\subsection{L'idea fondamentale}

La meccanica newtoniana descrive la gravità come una \emph{forza} che agisce
istantaneamente a distanza: ogni massa attrae ogni altra in proporzione al
prodotto delle masse e in ragione inversa del quadrato della distanza. Questa
descrizione funziona splendidamente nel regime delle velocità basse e dei campi
deboli, ma si incrina quando le velocità diventano relativistiche o i campi
gravitazionali estremamente intensi.

Einstein propose un'idea radicalmente diversa. La gravità non è una forza:
è una manifestazione della \emph{curvatura dello spaziotempo}. La presenza di
massa-energia deforma la struttura geometrica dello spazio e del tempo, e i
corpi \enquote{cadono} semplicemente seguendo le traiettorie più dritte possibili
(le \emph{geodetiche}) in questo spazio curvo. Come disse John Wheeler:

\begin{quote}
  \itshape
  ``La materia dice allo spaziotempo come curvarsi,
    lo spaziotempo dice alla materia come muoversi.''
\end{quote}

\subsection{Le equazioni di campo di Einstein}

Le equazioni di campo di Einstein, pubblicate nel novembre 1915, sono il cuore
della Relatività Generale. Sono dieci equazioni tensoriali (in realtà sei
indipendenti, per simmetrie) che legano la geometria dello spaziotempo alla
distribuzione di massa-energia:

\begin{physbox}{Equazioni di Campo di Einstein (1915)}
\begin{equation}
  \Gmunu + \Lambda \gmunu = \frac{8\pi G}{c^4}\, \Tmunu
  \label{eq:einstein}
\end{equation}
dove:
$\quad \Gmunu = \Rmunu - \tfrac{1}{2} R\, \gmunu$ \quad (tensore di Einstein)
\end{physbox}

\medskip
Ogni simbolo merita una spiegazione:

\begin{description}
  \item[$\gmunu$] Il \textbf{tensore metrico}: la quantità fondamentale che
    descrive la geometria dello spaziotempo. In spazio piatto (Minkowski) è
    $\mink = \text{diag}(-1,+1,+1,+1)$. In uno spaziotempo curvo $\gmunu$
    dipende dalle coordinate e dalla posizione. Definisce l'elemento di linea:
    \[
      ds^2 = \gmunu\, dx^\mu dx^\nu
    \]
    che generalizza il teorema di Pitagora a quattro dimensioni curve.

  \item[$\Rmunu$] Il \textbf{tensore di Ricci}: contrazione del tensore di
    Riemann $R^\rho_{\ \sigma\mu\nu}$, cioè $\Rmunu = R^\rho_{\ \mu\rho\nu}$.
    Il tensore di Riemann misura quanto un vettore ruota quando viene
    \emph{trasportato parallelamente} lungo una curva chiusa — una misura
    intrinseca della curvatura che non dipende dal sistema di coordinate.

  \item[$R$] La \textbf{scalare di Ricci}: $R = g^{\munu}\Rmunu$, la traccia
    del tensore di Ricci. Rappresenta la curvatura scalare media dello
    spaziotempo.

  \item[$\Lambda$] La \textbf{costante cosmologica}: introdotta da Einstein
    nel 1917 per ottenere un universo statico, poi abbandonata dopo la scoperta
    dell'espansione di Hubble (1929), e infine riscoperta come energia del vuoto
    responsabile dell'espansione accelerata dell'universo. Il suo valore
    misurato è $\Lambda \approx 1.1 \times 10^{-52}\ \text{m}^{-2}$.

  \item[$\Tmunu$] Il \textbf{tensore energia-impulso}: descrive la distribuzione
    di energia, pressione, densità e flusso di momento della materia. Per il
    vuoto $\Tmunu = 0$, e le equazioni di campo diventano $\Gmunu = 0$ (o
    equivalentemente $\Rmunu = 0$, dette \emph{equazioni di vuoto}).

  \item[$G$, $c$] La costante di gravitazione universale e la velocità della
    luce. Il fattore $8\pi G/c^4 \approx 2 \times 10^{-43}\ \text{N}^{-1}$
    è estremamente piccolo: la geometria dello spaziotempo risponde debolmente
    alla presenza di energia, il che spiega perché gli effetti relativistici
    siano trascurabili nella vita quotidiana.
\end{description}

\begin{curiosity}{Perché $8\pi$?}
  Il fattore $8\pi G/c^4$ non è arbitrario: è fissato richiedendo che nel
  limite di campo gravitazionale debole e velocità non relativistiche, le
  equazioni di Einstein si riducano alla legge di Poisson newtoniana
  $\nabla^2\Phi = 4\pi G\rho$. Questo \emph{limite newtoniano} è il punto di
  contatto tra le due teorie e vincola univocamente la costante.
\end{curiosity}

\subsection{I simboli di Christoffel e le geodetiche}

Dato un tensore metrico $\gmunu$, i \textbf{simboli di Christoffel} (o
\emph{connessione di Levi-Civita}) sono:

\begin{equation}
  \Gamma^\lambda_{\ \mu\nu}
  = \frac{1}{2}\, g^{\lambda\sigma}
    \left(\partial_\mu g_{\nu\sigma}
        + \partial_\nu g_{\mu\sigma}
        - \partial_\sigma g_{\mu\nu}\right)
  \label{eq:christoffel}
\end{equation}

Descrivono come variano le basi del sistema di coordinate nello spazio curvo.
Non sono un tensore (si trasformano in modo diverso) ma compaiono nell'equazione
delle \textbf{geodetiche} — le generalizzazioni relativistiche delle linee rette:

\begin{equation}
  \frac{d^2 x^\lambda}{d\tau^2}
  + \Gamma^\lambda_{\ \mu\nu}\frac{dx^\mu}{d\tau}\frac{dx^\nu}{d\tau} = 0
  \label{eq:geodetica}
\end{equation}

dove $\tau$ è il \emph{tempo proprio} dell'osservatore (il tempo misurato dal
suo orologio). Per la luce, che ha $ds^2 = 0$, si usa un parametro affine
$\lambda$ al posto di $\tau$. Queste equazioni — le leggi del moto in
Relatività Generale — sostituiscono $F = ma$ newtoniana.

\newpage

% ═════════════════════════════════════════════════════════════════════════════
\section{La Metrica di Schwarzschild}
% ═════════════════════════════════════════════════════════════════════════════

\subsection{La prima soluzione esatta}

Karl Schwarzschild trovò la prima soluzione esatta delle equazioni di Einstein
nel dicembre 1915, poche settimane dopo la pubblicazione delle equazioni stesse,
mentre si trovava al fronte orientale durante la Prima Guerra Mondiale. Morì
pochi mesi dopo per una malattia contratta in trincea.

La soluzione descrive lo spaziotempo \emph{esterno} a un oggetto sferico,
non rotante, nel vuoto ($\Tmunu = 0$). Il \textbf{teorema di Birkhoff} (1923)
dimostra che questa è la \emph{unica} soluzione a simmetria sferica nel vuoto:
qualunque distribuzione sferica di materia genera all'esterno esattamente la
metrica di Schwarzschild.

\subsection{L'elemento di linea}

\begin{physbox}{Metrica di Schwarzschild}
\begin{equation}
  ds^2 = -\!\left(1 - \frac{\RS}{r}\right)c^2\,dt^2
         + \left(1 - \frac{\RS}{r}\right)^{\!\!-1} dr^2
         + r^2\,d\theta^2 + r^2\sin^2\theta\, d\varphi^2
  \label{eq:schwarzschild}
\end{equation}
con il \textbf{raggio di Schwarzschild}:
\[
  \RS = \frac{2GM}{c^2}
\]
\end{physbox}

\medskip
In unità geometrizzate ($G = c = M = 1$), $\RS = 2$.

\subsection{Il significato fisico dei coefficienti}

La metrica \eqref{eq:schwarzschild} ha una struttura molto ricca:

\begin{description}
  \item[Coefficiente temporale $g_{tt} = -(1-\RS/r)$.]
    A grandi distanze ($r \to \infty$) vale $-1$: lo spazio è piatto e il
    tempo scorre normalmente. All'avvicinarsi a $r = \RS$, il coefficiente
    tende a zero: il \emph{tempo si ferma} per un osservatore lontano.
    Per $r < \RS$, $g_{tt}$ diventa positivo — tempo e spazio \emph{scambiano
    ruolo}: la coordinata $r$ diventa temporale e la caduta verso $r=0$
    è inevitabile come il passare del tempo.

  \item[Coefficiente radiale $g_{rr} = (1-\RS/r)^{-1}$.]
    Diverge per $r \to \RS$: le distanze radiali vengono dilatate
    indefinitamente. Tuttavia questa è solo una singolarità della coordinata,
    non fisica.

  \item[Parte angolare $r^2\,d\Omega^2$.]
    Identica a quella dello spazio piatto: la geometria delle sfere concentriche
    non viene alterata dalla curvatura radiale. Il raggio $r$ è in realtà la
    coordinata areale: una sfera a coordinata $r$ ha area $4\pi r^2$.
\end{description}

\subsection{Le due singolarità}

\subsubsection{$r = 0$: singolarità fisica}

All'origine, la curvatura diverge. In particolare l'invariante di Kretschner:
\begin{equation}
  K = R_{\mu\nu\rho\sigma}R^{\mu\nu\rho\sigma}
  = \frac{48 G^2 M^2}{c^4 r^6} \xrightarrow{r\to 0} \infty
\end{equation}
è un invariante scalare (non dipende dal sistema di coordinate), quindi questa
divergenza è reale. La Relatività Generale smette di essere valida qui: si
suppone che la meccanica quantistica — che la RG ignora completamente —
regolarizzi questa singolarità. Una teoria della \emph{gravità quantistica},
ancora mancante, è necessaria per descrivere $r = 0$.

\subsubsection{$r = \RS$: singolarità della coordinata (orizzonte degli eventi)}

Questa singolarità è un artificio del sistema di riferimento. Passando alle
\textbf{coordinate di Kruskal-Szekeres}:
\begin{align}
  T &= \left(\frac{r}{\RS} - 1\right)^{1/2} e^{r/(2\RS)} \sinh\!\left(\frac{ct}{2\RS}\right)\\
  X &= \left(\frac{r}{\RS} - 1\right)^{1/2} e^{r/(2\RS)} \cosh\!\left(\frac{ct}{2\RS}\right)
\end{align}
la metrica diventa regolare all'orizzonte. Fisicamente, $r = \RS$ è la
\textbf{superficie causale}: la frontiera da cui nessun segnale può sfuggire.

\subsubsection{Valori numerici del raggio di Schwarzschild}

\begin{table}[h!]
  \centering
  \begin{tabular}{llll}
    \toprule
    \textbf{Oggetto} & \textbf{Massa} & $\RS$ & \textbf{Densità media al collasso}\\
    \midrule
    Terra    & $5.97\times10^{24}$ kg & $8.9$ mm & $2\times10^{27}$ kg/m$^3$ \\
    Sole     & $1.99\times10^{30}$ kg & $2.95$ km & $1.8\times10^{19}$ kg/m$^3$ \\
    M87*     & $6.5\times10^9\, M_\odot$ & $\approx 1.9\times10^{10}$ km & $\approx 2$ kg/m$^3$ \\
    Sgr A*   & $4.15\times10^6\, M_\odot$ & $\approx 1.2\times10^7$ km & $\approx 5\times10^3$ kg/m$^3$ \\
    \bottomrule
  \end{tabular}
  \caption{Raggi di Schwarzschild per vari oggetti. Si noti che per i buchi
           neri supermassicci la densità media è straordinariamente bassa.}
\end{table}

\begin{curiosity}{M87* ha densità inferiore all'aria!}
  Il buco nero M87*, con $\RS \approx 1.9 \times 10^{10}$ km, ha una densità
  media al collasso (calcolata come massa su volume della sfera di raggio $\RS$)
  di circa $2$ kg/m$^3$ — inferiore alla densità dell'aria a livello del mare
  ($\approx 1.2$ kg/m$^3$). Questo sembra paradossale ma è corretto: per
  masse sufficientemente grandi, il collasso avviene a densità sempre più
  basse, perché $\RS \propto M$ mentre il volume scala come $M^3$.
  La sensazione di caduta vertiginosa all'orizzonte dipende dall'accelerazione
  di marea $\sim GM/\RS^3 \propto 1/M^2$: per M87* è praticamente zero
  all'orizzonte. Un astronauta non sentirebbe nulla attraversandolo.
\end{curiosity}

\newpage

% ═════════════════════════════════════════════════════════════════════════════
\section{La Metrica di Kerr e la Rotazione}
% ═════════════════════════════════════════════════════════════════════════════

\subsection{La soluzione di Kerr}

Nel 1963, quarantott'anni dopo Schwarzschild, Roy Kerr trovò la soluzione per
un buco nero \emph{in rotazione} con momento angolare $J$. Si tratta della
soluzione più realistica per i buchi neri astrofisici, poiché tutti i corpi
celesti hanno momento angolare non nullo per conservazione angolare durante
il collasso.

Si introduce il \textbf{parametro di spin} $a = J/(Mc)$, con le dimensioni
di una lunghezza, e il parametro adimensionale $a^* = ac/(GM) \in [0, 1]$.
Due funzioni ausiliarie semplificano la scrittura:
\begin{align}
  \Sigma &= r^2 + a^2\cos^2\theta \\
  \Delta &= r^2 - \RS\, r + a^2
\end{align}

\begin{physbox}{Metrica di Kerr (coordinate di Boyer-Lindquist)}
\begin{equation}
  ds^2 = -\frac{\Delta}{\Sigma}\!\left(dt - a\sin^2\!\theta\, d\varphi\right)^2
        + \frac{\Sigma}{\Delta}\,dr^2
        + \Sigma\, d\theta^2
        + \frac{\sin^2\!\theta}{\Sigma}\!\left(a\,dt - (r^2+a^2)d\varphi\right)^2
  \label{eq:kerr}
\end{equation}
\end{physbox}

\subsection{I due orizzonti}

Le radici di $\Delta = 0$ danno i due orizzonti:
\begin{equation}
  r_\pm = \frac{\RS}{2} \pm \sqrt{\left(\frac{\RS}{2}\right)^2 - a^2}
\end{equation}
Per $a = 0$ si recupera $r_+ = \RS$ (Schwarzschild). Per $a = GM/c$
(\textbf{buco nero estremale}), $r_+ = r_- = M$: i due orizzonti coincidono
e la struttura causale del buco nero cambia radicalmente.

\subsection{L'ergosfera e il frame dragging}

Per la metrica di Kerr, la superficie $g_{tt} = 0$ non coincide con l'orizzonte
ma si trova a:
\begin{equation}
  r_{\text{ergo}}(\theta) = \frac{\RS}{2} + \sqrt{\left(\frac{\RS}{2}\right)^2 - a^2\cos^2\theta}
\end{equation}
La regione $r_+ < r < r_{\text{ergo}}$ è l'\textbf{ergosfera}. Qui, il
coefficiente $g_{tt}$ cambia segno: qualunque osservatore è costretto a
co-ruotare con il buco nero (\emph{frame dragging}), anche se impiega propulsori
per restare fermo. Non può esistere nessun osservatore a riposo rispetto alle
stelle lontane dentro l'ergosfera.

\begin{curiosity}{Il Processo di Penrose}
  Roger Penrose dimostrò nel 1969 che è possibile estrarre energia rotazionale
  da un buco nero di Kerr tramite un processo che avviene nell'ergosfera.
  Se una particella si divide in due dentro l'ergosfera, una delle due può
  avere energia \emph{negativa} (permessa perché $g_{tt} > 0$) e cadere
  nel buco nero, mentre l'altra esce con energia superiore a quella della
  particella originale. Il buco nero rallenta. L'energia massima estraibile
  è $1 - 1/\sqrt{2} \approx 29\%$ della massa del buco nero estremale.
  I jet relativistici osservati in molti quasar potrebbero sfruttare un
  meccanismo analogo (processo di Blandford-Znajek).
\end{curiosity}

\newpage

% ═════════════════════════════════════════════════════════════════════════════
\section{Le Geodetiche Nulle e il Moto della Luce}
% ═════════════════════════════════════════════════════════════════════════════

\subsection{Le costanti del moto}

La metrica di Schwarzschild è stazionaria (non dipende da $t$) e a simmetria
azimutale (non dipende da $\varphi$). Per il teorema di Noether, esistono due
costanti del moto per ogni fotone:

\begin{equation}
  E = -g_{tt}\frac{dt}{d\lambda} = \left(1 - \frac{\RS}{r}\right)\frac{dt}{d\lambda}
  \qquad \text{(energia)}
\end{equation}
\begin{equation}
  L = g_{\varphi\varphi}\frac{d\varphi}{d\lambda} = r^2\frac{d\varphi}{d\lambda}
  \qquad \text{(momento angolare)}
\end{equation}

Il \textbf{parametro d'impatto} $b = L/E$ ha un significato geometrico chiaro:
è la distanza perpendicolare tra il raggio e il centro del buco nero misurata
lontano da esso. Un raggio radiale ha $b = 0$, un raggio tangente lontano ha
$b \to \infty$.

\subsection{Il potenziale efficace per la luce}

Il vincolo $ds^2 = 0$ per la geodetica nulla nel piano equatoriale
($\theta = \pi/2$) dà:

\begin{equation}
  \left(\frac{dr}{d\lambda}\right)^2 = E^2 - \frac{L^2}{r^2}\left(1 - \frac{\RS}{r}\right)
  \equiv E^2 - V_{\text{eff}}(r)
\end{equation}

Il \textbf{potenziale efficace} per la luce è:
\begin{equation}
  V_{\text{eff}}(r) = \frac{L^2}{r^2}\left(1 - \frac{\RS}{r}\right)
  = \frac{L^2}{r^2} - \frac{L^2 \RS}{r^3}
  \label{eq:veff}
\end{equation}

\begin{figure}[h!]
  \centering
  \begin{tikzpicture}
    \begin{axis}[
      width=10cm, height=6cm,
      xlabel={$r/\RS$}, ylabel={$V_{\text{eff}}/L^2$},
      xmin=1, xmax=10, ymin=-0.01, ymax=0.055,
      grid=major, grid style={dashed, gray!30},
      axis lines=left,
      tick label style={font=\small},
      label style={font=\small\color{NavyBlue}},
      every axis plot/.append style={thick}
    ]
      \addplot[NavyBlue, domain=1.1:10, samples=200]
        { 1/x^2 * (1 - 2/x) };
      \addplot[CrimsonRed, dashed, domain=1.1:10]
        { 1/(27) };
      \addplot[ForestGreen, dotted, domain=1.5:1.5, ymin=0, ymax=0.055]
        coordinates {(1.5, 0) (1.5, 1/(27))};
      \node[font=\small, color=CrimsonRed] at (axis cs:8, 0.040) {$V_{\max} = \frac{1}{27}$};
      \node[font=\small, color=ForestGreen] at (axis cs:1.8, 0.005) {$r=1.5\,\RS$};
    \end{axis}
  \end{tikzpicture}
  \caption{Potenziale efficace per la luce nella metrica di Schwarzschild
           (in unità $L^2 = 1$). Il massimo a $r = 1.5\,\RS$ corrisponde
           alla fotosfera. I raggi con $b^2 = L^2/E^2 < 27M^2$ cadono
           nel buco nero; quelli con $b^2 > 27M^2$ sfuggono.}
\end{figure}

\subsection{La fotosfera e il parametro d'impatto critico}

Il massimo di $V_{\text{eff}}$ si trova ponendo $dV_{\text{eff}}/dr = 0$:

\begin{equation}
  \frac{dV_{\text{eff}}}{dr} = -\frac{2L^2}{r^3} + \frac{3L^2\RS}{r^4} = 0
  \quad\Longrightarrow\quad
  r_{\text{ph}} = \frac{3\RS}{2} = \frac{3GM}{c^2}
\end{equation}

Questo è il raggio della \textbf{fotosfera} (o \emph{photon sphere}): l'orbita
circolare instabile dei fotoni. Un fotone con esattamente il parametro
d'impatto critico:

\begin{physbox}{Parametro d'impatto critico}
\begin{equation}
  \bc = \sqrt{27}\, M = \frac{3\sqrt{3}}{2}\, \RS \approx 2.598\,\RS
  \label{eq:bcrit}
\end{equation}
\end{physbox}

orbita indefinitamente sulla fotosfera. Questo valore determina la dimensione
apparente dell'\textbf{ombra del buco nero}: tutti i raggi con $b < \bc$ cadono
nell'orizzonte; quelli con $b > \bc$ sfuggono. Per un osservatore lontano,
l'ombra ha raggio angolare:
\begin{equation}
  \alpha_{\text{shadow}} = \frac{\bc}{D} = \frac{3\sqrt{3}\,GM}{c^2 D}
\end{equation}
dove $D$ è la distanza. Per M87* ($D = 16.8$ Mpc, $M = 6.5 \times 10^9 M_\odot$):
$\alpha_{\text{shadow}} \approx 20\ \mu\text{as}$ — confermato entro il 17\%
dalla fotografia EHT del 2019.

\begin{curiosity}{Perché l'ombra è più grande dell'orizzonte?}
  L'ombra del buco nero ha raggio apparente $\bc \approx 2.6\,\RS$, molto
  maggiore dell'orizzonte degli eventi a $\RS$. Il motivo è il lensing
  gravitazionale: i raggi che passano \emph{vicino} all'orizzonte ma non
  lo attraversano vengono deflessi così tanto che, visti dall'osservatore,
  sembrano provenire dall'interno dell'ombra. L'ombra è dunque la \emph{proiezione
  della fotosfera}, non dell'orizzonte.
\end{curiosity}

\subsection{L'equazione di Binet relativistica}

Sostituendo $u = 1/r$ (variabile di Binet), la traiettoria dei fotoni diventa:

\begin{physbox}{Equazione di Binet per la luce (RG)}
\begin{equation}
  \frac{d^2 u}{d\varphi^2} = -u + \frac{3\RS}{2}\, u^2
  \label{eq:binet}
\end{equation}
\end{physbox}

Il termine $-u$ è la traiettoria rettilinea in spazio piatto (con $b = 1/u_0$
costante). Il termine $\frac{3\RS}{2}\, u^2$ è la \textbf{correzione relativistica}:
è questo unico termine a generare tutto il lensing gravitazionale, la fotosfera,
l'ombra e il photon ring.

Riscriviamo come sistema del primo ordine per l'integrazione numerica:
\begin{align}
  \frac{du}{d\varphi} &= w \\
  \frac{dw}{d\varphi} &= -u + \frac{3\RS}{2}\, u^2
  \label{eq:ode_system}
\end{align}
Questo è il sistema che integriamo con Runge-Kutta nel codice GLSL.

\newpage

% ═════════════════════════════════════════════════════════════════════════════
\section{Il Disco di Accrescimento}
% ═════════════════════════════════════════════════════════════════════════════

\subsection{Perché un disco?}

Il materiale che cade verso un buco nero porta con sé momento angolare.
Per conservazione del momento angolare, non può cadere direttamente verso
il centro ma forma una struttura con simmetria azimutale: il \textbf{disco
di accrescimento}. Il gas nel disco perde lentamente energia e momento angolare
per attrito viscoso (in realtà instabilità magnetoidrodinamiche — l'instabilità
di Magneto-Rotazionale, o MRI), spiralando lentamente verso l'interno.

\subsection{L'orbita circolare stabile più interna (ISCO)}

Non tutte le orbite circolari sono stabili. Nella metrica di Schwarzschild,
l'orbita circolare più interna stabile — detta \textbf{ISCO} dall'inglese
\textit{Innermost Stable Circular Orbit} — si trova a:

\begin{equation}
  \Risco = 3\,\RS = \frac{6GM}{c^2}
  \label{eq:isco}
\end{equation}

Il materiale dentro l'ISCO perde stabilità orbitale e cade rapidamente verso
l'orizzonte senza dissipare ulteriore energia — il disco finisce a $\Risco$.
Per la metrica di Kerr, l'ISCO dipende dallo spin:
\begin{itemize}
  \item Spin zero (Schwarzschild): $\Risco = 6M$
  \item Spin massimo pro-grado: $\Risco = M$ (il bordo del buco nero stesso)
  \item Spin massimo retrogrado: $\Risco = 9M$
\end{itemize}
Questo significa che i buchi neri in rapida rotazione accrescono in modo
molto più efficiente: il materiale può spiralare molto più vicino e rilasciare
molta più energia gravitazionale prima di cadere.

\subsection{Il profilo di temperatura: modello di Novikov-Thorne}

Il modello di Novikov-Thorne (1973) descrive la struttura termica di un disco
geometricamente sottile e otticamente spesso. Il flusso di energia radiata
per unità di superficie del disco è:

\begin{equation}
  Q(r) = \frac{3GM\dot{M}}{8\pi r^3}
         \left(1 - \sqrt{\frac{\Risco}{r}}\right)
         \cdot f(r)
\end{equation}

dove $\dot{M}$ è il tasso di accrescimento (massa per unità di tempo) e $f(r)$
è un fattore di correzione relativistico. Nello shader usiamo l'approssimazione
di corpo nero locale $Q = \sigma_{\text{SB}} T^4$, ottenendo il profilo di
temperatura approssimato:

\begin{physbox}{Profilo di temperatura del disco}
\begin{equation}
  T(r) \propto \left(\frac{r}{\Risco}\right)^{-3/4}
               \left(1 - \sqrt{\frac{\Risco}{r}}\right)^{1/4}
  \label{eq:temperature_profile}
\end{equation}
\end{physbox}

\begin{description}
  \item[Il fattore $r^{-3/4}$] proviene dal flusso kepleriano: zone più interne
    sono più calde perché l'energia potenziale gravitazionale rilasciata per
    unità di area aumenta verso l'interno.
  \item[Il fattore $(1 - \sqrt{\Risco/r})^{1/4}$] impone $T \to 0$ all'ISCO:
    il bordo interno del disco è freddo perché il materiale scivola dentro
    senza dissipare ulteriore calore.
\end{description}

Le temperature raggiunte nel disco variano da:
\begin{itemize}
  \item $\sim 10^4$ K al bordo esterno (emissione nell'infrarosso e nel visibile)
  \item $\sim 10^6$–$10^7$ K vicino all'ISCO per buchi neri stellari (emissione in raggi X)
  \item $\sim 10^5$ K per i buchi neri supermassicci (UV e raggi X molli)
\end{itemize}

\subsection{L'efficienza di accrescimento}

L'energia rilasciata per unità di massa in caduta libera newtoniana sarebbe
$GM/(2r)$ — senza limite teorico. Con il profilo di Novikov-Thorne, l'efficienza
di conversione massa-energia è:

\begin{equation}
  \eta_{\text{Schwarzschild}} = 1 - \sqrt{1 - \frac{2}{3}\cdot\frac{\RS}{\Risco}}
  = 1 - \sqrt{\frac{8}{9}} \approx 5.7\%
\end{equation}

Per confronto, la fusione nucleare (idrogeno in elio) raggiunge solo $\sim 0.7\%$.
I buchi neri in accrescimento sono i motori più efficienti dell'universo.
Per lo spin massimo di Kerr, $\eta$ sale fino a $\approx 42\%$.

\newpage

% ═════════════════════════════════════════════════════════════════════════════
\section{Effetti Osservativi Relativistici}
% ═════════════════════════════════════════════════════════════════════════════

\subsection{Il redshift gravitazionale}

Un fotone emesso a raggio $r$ con frequenza $\nu_e$ e osservato lontano dal
campo gravitazionale arriva con frequenza:

\begin{physbox}{Redshift Gravitazionale di Schwarzschild}
\begin{equation}
  \nu_\infty = \nu_e \sqrt{1 - \frac{\RS}{r}}
  \label{eq:redshift}
\end{equation}
\end{physbox}

Per la legge di Wien ($\nu_{\max} \propto T$), la temperatura apparente
del disco viene ridotta dello stesso fattore:
\begin{equation}
  \Tobs(r) = \Tem(r) \cdot \sqrt{1 - \frac{\RS}{r}}
\end{equation}

All'orizzonte ($r \to \RS$): $\nu_\infty \to 0$. Qualunque cosa accada
all'orizzonte è invisibile per noi: i fotoni perdono tutta la loro energia
in frequenza infinitamente più bassa. Questo è il motivo per cui l'ombra
del buco nero è nera.

\begin{curiosity}{Il redshift gravitazionale misurato sulla Terra}
  Il redshift gravitazionale è misurabile anche sulla Terra. Un orologio
  al piano terra di un grattacielo scorre più lentamente di uno al piano
  superiore. Nel 2010, con orologi ottici al literbio, è stato misurato
  un $\Delta\nu/\nu \approx 4 \times 10^{-17}$ su una differenza di quota
  di soli 33 cm — la previsione di Einstein confermata con precisione
  al livello di $10^{-4}$.
\end{curiosity}

\subsection{L'effetto Doppler relativistico}

Il materiale nel disco ruota a velocità relativistiche. La velocità orbitale
kepleriana nella metrica di Schwarzschild per un'orbita circolare a raggio $r$ è:

\begin{equation}
  v_K = \sqrt{\frac{GM}{r}} \cdot \frac{1}{\sqrt{1 - \RS/r}} \cdot \frac{1}{c}
\end{equation}

(in unità dove $c = 1$). All'ISCO ($r = 3\RS$):
\begin{equation}
  v_K(\Risco) = \sqrt{\frac{M}{6M}} \cdot \frac{1}{\sqrt{1-1/3}} = \frac{1}{\sqrt{6}} \cdot \frac{\sqrt{3}}{\sqrt{2}} = \frac{1}{2}c
\end{equation}
cioè metà della velocità della luce — profondamente relativistico.

Il fattore Doppler per un osservatore che vede il materiale muoversi con
velocità $\beta = v/c$ a un angolo $\psi$ rispetto alla linea di vista è:

\begin{physbox}{Fattore Doppler Relativistico}
\begin{equation}
  \mathcal{D} = \frac{1}{\gamma\,(1 - \beta\cos\psi)},
  \qquad \gamma = \frac{1}{\sqrt{1-\beta^2}}
  \label{eq:doppler}
\end{equation}
\end{physbox}

Il flusso osservato (potenza per unità di area e frequenza) scala come:
\begin{equation}
  F_{\text{obs}} = \mathcal{D}^3 \cdot F_{\text{em}}
\end{equation}
Il cubo compare perché:
\begin{enumerate}
  \item un fattore $\mathcal{D}$ proviene dalla variazione dell'energia dei fotoni
        ($E_{\text{obs}} = \mathcal{D} E_{\text{em}}$),
  \item un fattore $\mathcal{D}$ proviene dall'aberrazione (la sezione trasversale
        apparente del fascio varia),
  \item un fattore $\mathcal{D}$ proviene dalla frequenza di arrivo dei fotoni
        (rate = $\mathcal{D}$ volte il rate di emissione).
\end{enumerate}

\subsection{L'asimmetria luminosa del disco}

La combinazione di redshift gravitazionale e Doppler produce un'asimmetria
caratteristica nell'immagine del disco:
\begin{itemize}
  \item Il lato in avvicinamento ($\cos\psi > 0$, $\mathcal{D} > 1$): più
        brillante, più blu, spostato verso l'esterno per il lensing.
  \item Il lato in allontanamento ($\cos\psi < 0$, $\mathcal{D} < 1$): più
        scuro, più rosso, compresso verso l'interno.
\end{itemize}

Questa asimmetria è chiaramente visibile nella fotografia di M87* del 2019:
il lato sinistro (inferiore in realtà, per la geometria di quella galassia)
appare circa 3 volte più luminoso. Il fattore $\mathcal{D}^3$ può variare
da $\sim 0.1$ a $\sim 3$-$4$ lungo il disco.

\subsection{Il photon ring}

I fotoni con parametro d'impatto $b \approx \bc$ compiono quasi un'orbita
completa attorno alla fotosfera prima di sfuggire. Questi fotoni formano un
\textbf{photon ring}: un anello brillante e sottilissimo all'interno del quale
si trova l'ombra del buco nero. In realtà il photon ring è composto da
infinite sub-strutture:
\begin{itemize}
  \item Anello di ordine 0: luce che arriva direttamente dal disco senza
        orbite aggiuntive (la parte principale del disco).
  \item Anello di ordine 1: luce che ha orbitato mezza volta attorno al BH.
  \item Anello di ordine $n$: luce che ha orbitato $n/2$ volte.
\end{itemize}
La luminosità decresce geometricamente: ogni anello successivo è circa
$e^{-\pi}$ più debole del precedente. Gli anelli di ordine $\geq 2$ sono
praticamente inosservabili con le tecnologie attuali.

\begin{curiosity}{La vera risoluzione di EHT}
  L'Event Horizon Telescope (EHT) è un interferometro radio con baseline pari
  al diametro terrestre ($\sim 10^4$ km). Operando a $\lambda = 1.3$ mm,
  raggiunge una risoluzione angolare di $\theta \sim \lambda/D \approx
  20\ \mu$as (microarcosecondi). Per confronto, l'occhio umano risolve
  circa $10^8$ volte peggio ($\sim 1'$ d'arco). La fotografia di M87* del
  2019 mostra una struttura di circa 40 $\mu$as di diametro — equivalente
  a vedere una pallina da golf a distanza della Luna.
\end{curiosity}

\newpage

% ═════════════════════════════════════════════════════════════════════════════
\section{Il Redshift Gravitazionale e la Radiazione di Hawking}
% ═════════════════════════════════════════════════════════════════════════════

\subsection{La temperatura di Hawking}

Nel 1974 Stephen Hawking dimostrò, combinando Relatività Generale e meccanica
quantistica dei campi, che i buchi neri non sono completamente neri: emettono
radiazione termica con \textbf{temperatura di Hawking}:

\begin{physbox}{Temperatura di Hawking}
\begin{equation}
  T_H = \frac{\hbar c^3}{8\pi G M k_B}
      \approx \frac{6.2 \times 10^{-8}\ \text{K}}{M/M_\odot}
\end{equation}
\end{physbox}

Per un buco nero di massa solare, $T_H \approx 62$ nK — completamente
irrilevabile rispetto alla temperatura del fondo cosmico a microonde
($T_{\text{CMB}} \approx 2.7$ K). I buchi neri stellari sono di fatto
eterni su scale cosmologiche.

Al contrario, un ipotetico micro buco nero di massa $10^{12}$ kg
(un meteorite di 1 km) avrebbe $T_H \sim 10^{12}$ K e evapora
nell'ordine del tempo dell'universo.

\subsection{Il paradosso dell'informazione}

Se un buco nero evapora completamente per radiazione di Hawking, le
informazioni del materiale caduto dentro — in linea di principio
registrate nello stato quantistico di ogni particella — dove vanno?
La meccanica quantistica richiede l'unitarietà dell'evoluzione temporale:
l'informazione non può essere distrutta. La Relatività Generale sembra
invece suggerire che l'informazione sia perduta irrecuperabilmente.

Questo \textbf{paradosso dell'informazione} (Hawking, 1976) è uno dei
problemi aperti più profondi della fisica teorica. Recenti progressi
suggeriscono che l'informazione emerga nell'ultima fase dell'evaporazione
attraverso correlazioni quantistiche nella radiazione di Hawking (isole
di entanglement, pagine di entanglement). La risoluzione completa richiede
probabilmente una teoria della gravità quantistica.

\newpage

% ═════════════════════════════════════════════════════════════════════════════
%  PARTE II — IMPLEMENTAZIONE
% ═════════════════════════════════════════════════════════════════════════════

\section{Il Ray Tracing Relativistico}

\subsection{Filosofia dell'algoritmo}

Renderizzare un buco nero richiede di tracciare il percorso dei fotoni in uno
spaziotempo curvo. A differenza del ray tracing classico (dove la luce viaggia
in linea retta), qui i raggi si curvano seguendo le geodetiche della metrica
di Schwarzschild.

La strategia è il \textbf{ray tracing all'indietro nel tempo} (\emph{backward
ray tracing}): per ogni pixel dell'immagine, si lancia un raggio dal punto
di vista dell'osservatore e si integra la geodetica \emph{all'indietro} verso
la sorgente. Questo è fisicamente legittimo per la time-reversal symmetry
delle equazioni di Maxwell nel vuoto.

\subsection{Il metodo di Runge-Kutta al quarto ordine}

L'equazione di Binet \eqref{eq:binet} è un'equazione differenziale ordinaria
del secondo ordine. Riscritta come sistema del primo ordine
$(u, w) = (1/r, du/d\varphi)$:

\begin{equation}
  \frac{d}{d\varphi}\begin{pmatrix}u \\ w\end{pmatrix}
  = \mathbf{F}(u, w)
  = \begin{pmatrix} w \\ -u + \tfrac{3\RS}{2}\,u^2 \end{pmatrix}
\end{equation}

Il metodo \textbf{Runge-Kutta al quarto ordine} (RK4) approssima la soluzione
con errore locale $O(h^5)$ e globale $O(h^4)$:

\begin{physbox}{Schema Runge-Kutta 4 (RK4)}
\begin{align}
  \mathbf{k}_1 &= \mathbf{F}(\mathbf{y}_n) \notag\\
  \mathbf{k}_2 &= \mathbf{F}\!\left(\mathbf{y}_n + \tfrac{h}{2}\mathbf{k}_1\right) \notag\\
  \mathbf{k}_3 &= \mathbf{F}\!\left(\mathbf{y}_n + \tfrac{h}{2}\mathbf{k}_2\right) \notag\\
  \mathbf{k}_4 &= \mathbf{F}\!\left(\mathbf{y}_n + h\,\mathbf{k}_3\right) \notag\\
  \mathbf{y}_{n+1} &= \mathbf{y}_n
    + \frac{h}{6}\left(\mathbf{k}_1 + 2\mathbf{k}_2 + 2\mathbf{k}_3 + \mathbf{k}_4\right)
  \label{eq:rk4}
\end{align}
\end{physbox}

RK4 valuta la funzione $\mathbf{F}$ quattro volte per passo: ai bordi ($k_1$,
$k_4$) e nel mezzo ($k_2$, $k_3$). La media pesata $1{:}2{:}2{:}1$ garantisce
la simmetria del metodo e la convergenza al quarto ordine.

\subsection{Condizioni di terminazione}

Durante l'integrazione, il raggio può:

\begin{enumerate}
  \item \textbf{Cadere nell'orizzonte} ($r \leq \RS$): il pixel è nero.
  \item \textbf{Intersecare il disco} (piano $y=0$, $\Risco \leq r \leq r_{\text{out}}$):
        si calcola il colore del disco con il profilo di Novikov-Thorne.
  \item \textbf{Sfuggire all'infinito} ($r > r_{\max}$): si recupera il
        campo stellare di sfondo, deflesso dal lensing.
\end{enumerate}

\subsection{Il colore finale del pixel}

La catena completa per ogni pixel è:

\begin{equation}
  \text{pixel} \xrightarrow{\text{ray cast}}
  (r_0, \mathbf{d}) \xrightarrow{\text{RK4, 512 passi}}
  \text{esito} \xrightarrow{\text{colore}}
  \begin{cases}
    \mathbf{0} & \text{orizzonte}\\
    \text{HSV}(T_{\text{obs}}, \mathcal{D}) & \text{disco}\\
    \text{stelle}(\mathbf{d}_{\text{deflessa}}) & \text{fuga}
  \end{cases}
\end{equation}

\subsection{Il tone mapping e la correzione gamma}

I calcoli fisici producono valori di radianza in spazio lineare e potenzialmente
illimitati. Due passaggi di post-processing adattano l'output al display:

\subsubsection{Tone mapping di Reinhard}

Comprime l'intervallo $[0, +\infty) \to [0, 1)$ preservando i dettagli nelle
alte luci:
\begin{equation}
  L_{\text{display}} = \frac{L_{\text{lineare}}}{L_{\text{lineare}} + 1}
\end{equation}

\subsubsection{Correzione gamma}

I monitor CRT (e i display moderni che ne ereditano la convenzione) hanno
una risposta non lineare. La correzione sRGB usa $\gamma = 2.2$:
\begin{equation}
  \text{sRGB} = L_{\text{display}}^{1/2.2}
\end{equation}
Senza questa correzione, le immagini calcolate in spazio lineare appaiono
troppo scure nelle mezze tonalità.

\newpage

% ═════════════════════════════════════════════════════════════════════════════
\section{Il Codice GLSL Commentato}
% ═════════════════════════════════════════════════════════════════════════════

\subsection{Struttura generale}

Lo shader è un \textbf{fragment shader} GLSL 3.30 che gira sulla GPU per ogni
pixel dello schermo in parallelo. Le variabili \texttt{uniform} sono passate
dalla CPU ogni frame:

\begin{itemize}
  \item \texttt{iResolution}: dimensioni in pixel dello schermo
  \item \texttt{iTime}: secondi dall'avvio (per l'animazione)
  \item \texttt{uInclination}: angolo di inclinazione della camera
\end{itemize}

\subsection{Funzioni chiave}

\subsubsection{geodesicRHS — La destra del sistema ODE}

Implementa $\mathbf{F}(u,w) = (w,\ -u + 1.5\,\RS\,u^2)$:

\begin{tcolorbox}[codebox, title={geodesicRHS.glsl}]
\begin{verbatim}
vec2 geodesicRHS(vec2 y) {
    float u    = y.x;     // u = 1/r
    float duPhi = y.y;    // du/dphi
    return vec2(duPhi, -u + 1.5 * RS * u * u);
    //          ^^^^^   ^^^^^^^^^^^^^^^^^^^^^^^^^^
    //        dy[0]/dphi   dy[1]/dphi (eq. di Binet)
}
\end{verbatim}
\end{tcolorbox}

\subsubsection{rk4Step — Un passo di Runge-Kutta 4}

\begin{tcolorbox}[codebox, title={rk4Step.glsl}]
\begin{verbatim}
vec2 rk4Step(vec2 y, float h) {
    vec2 k1 = geodesicRHS(y);
    vec2 k2 = geodesicRHS(y + 0.5*h*k1);
    vec2 k3 = geodesicRHS(y + 0.5*h*k2);
    vec2 k4 = geodesicRHS(y + h*k3);
    return y + (h/6.0) * (k1 + 2.0*k2 + 2.0*k3 + k4);
}
\end{verbatim}
\end{tcolorbox}

\subsubsection{accretionDiskColor — Colore del disco}

Implementa la catena:
$r \to T_{\text{em}}(r) \to T_{\text{obs}} \times \sqrt{1-\RS/r} \to
T_{\text{obs}} \times \mathcal{D} \to \text{blackbody} \times \mathcal{D}^3$

\begin{tcolorbox}[codebox, title={accretionDiskColor.glsl}]
\begin{verbatim}
// Profilo Novikov-Thorne (approssimato)
float x        = r / R_ISCO;
float tProfile = pow(x, -0.75) * pow(max(1.0 - 1.0/sqrt(x), 0.0), 0.25);
float T_emit   = 8000.0 * tProfile;       // Kelvin (normalizzati)

// Redshift gravitazionale: T_obs = T_emit * sqrt(1 - rs/r)
float gravRedshift = sqrt(max(1.0 - RS/r, 0.0));
float T_obs        = T_emit * gravRedshift;

// Velocità kepleriana GR e fattore Doppler
float beta  = clamp(sqrt(M/r) / sqrt(1.0 - RS/r), 0.0, 0.95);
float gamma = 1.0 / sqrt(1.0 - beta*beta);
float D     = 1.0 / (gamma * (1.0 - beta * sin(phi_angle)));

// Temperatura osservata (include Doppler)
T_obs *= D;

// Flusso: D^3 * redshift^4 * profilo
float flux = pow(D, 3.0) * pow(gravRedshift, 4.0) * tProfile;
\end{verbatim}
\end{tcolorbox}

\subsection{Complessità computazionale}

Per ogni pixel (tipicamente $1920 \times 1080 \approx 2 \times 10^6$ pixel):
\begin{itemize}
  \item 512 passi RK4, ciascuno con 4 valutazioni di \texttt{geodesicRHS}
  \item Total: $512 \times 4 = 2048$ valutazioni per pixel
  \item Total screen: $\sim 4 \times 10^9$ operazioni floating-point per frame
\end{itemize}

Una moderna GPU (es. NVIDIA RTX 4090, $\sim 82$ TFLOPS) gestisce questo in
circa $50$ ms, ottenendo $\sim 20$ fps. Riducendo i passi a 256 si raggiungono
60 fps con minima perdita di qualità.

\newpage

% ═════════════════════════════════════════════════════════════════════════════
\section{Formulazione Matematica Completa dello Shader}
% ═════════════════════════════════════════════════════════════════════════════

Riassumendo tutte le equazioni in una catena causale unica, il colore di un
pixel $(p_x, p_y)$ al tempo $t$ è:

\subsection*{Passo 1: Setup del raggio}

\begin{equation}
  \mathbf{d} = \left(\frac{p_x - W/2}{H},\ \frac{p_y - H/2}{H},\ 0.7\right)_{\!\text{norm}}
\end{equation}

\subsection*{Passo 2: Parametro d'impatto}

\begin{equation}
  b = r_0 \cdot |\mathbf{d} \times \hat{\mathbf{r}}_0|
\end{equation}

\subsection*{Passo 3: Integrazione RK4 dell'equazione di Binet}

\begin{equation}
  (u_{n+1}, w_{n+1}) = \text{RK4}\!\left[(u_n, w_n),\ h = \frac{\Delta\varphi}{r_n}\right]
  \quad \text{per } n = 0, 1, \dots, 511
\end{equation}

\subsection*{Passo 4: Colore}

\[
  \text{pixel} = \begin{cases}
    \mathbf{0}
      & r \leq \RS \\[4pt]
    \text{BB}\!\left[
      T_{\text{em}}(r) \cdot \sqrt{1-\RS/r} \cdot \mathcal{D}(\beta, \psi)
    \right] \cdot \mathcal{D}^3 \cdot (1-\RS/r)^2
      & r_{\text{ISCO}} \leq r \leq r_{\text{out}} \\[4pt]
    \text{stelle}(\mathbf{d}_{\text{lens}})
      + \text{photon ring}(|b - \bc|)
      & \text{altrimenti}
  \end{cases}
\]

dove $\text{BB}[T]$ è il colore di un corpo nero alla temperatura $T$,
approssimato dalla funzione \texttt{blackbodyColor} dello shader.

\subsection*{Passo 5: Post-processing}

\begin{equation}
  \text{output} = \left(\frac{\text{pixel}}{\text{pixel} + \mathbf{1}}\right)^{1/2.2}
\end{equation}

\newpage

% ═════════════════════════════════════════════════════════════════════════════
\section{Approfondimenti e Letture Consigliate}
% ═════════════════════════════════════════════════════════════════════════════

\subsection{Testi di riferimento}

\begin{description}
  \item[Misner, Thorne, Wheeler — \textit{Gravitation} (1973)]
    Il testo fondamentale della Relatività Generale, monumentale e completo.
    Il capitolo 25 tratta le geodetiche di Schwarzschild.

  \item[Wald — \textit{General Relativity} (1984)]
    Trattazione matematicamente rigorosa, ideale per chi ha basi di geometria
    differenziale. Utilizza la notazione con forme differenziali.

  \item[Carroll — \textit{Spacetime and Geometry} (2004)]
    Testo moderno e accessibile, con eccellente introduzione alla geometria
    differenziale e ai buchi neri.

  \item[Novikov, Thorne — in DeWitt \& DeWitt (1973)]
    L'articolo originale sul profilo di temperatura del disco di accrescimento.

  \item[Luminet — \textit{Image of a Spherical Black Hole} (1979)]
    La prima simulazione visiva di un buco nero, precursore storico degli
    shader moderni. L'autore calcolò l'immagine su carta millimetrata.

  \item[James et al. — \textit{Gravitational lensing in Interstellar} (2015)]
    Come il team di Kip Thorne realizzò il buco nero di Interstellar.
    Descrive il codice DNGR (Double Negative Gravitational Renderer),
    che risolve le geodetiche in 4D complete.
\end{description}

\subsection{Estensioni future dello shader}

La simulazione presentata in questo documento può essere estesa in vari modi:

\begin{enumerate}
  \item \textbf{Metrica di Kerr completa}: integrare le geodetiche 4D nella
    metrica di Kerr richiede l'uso della quarta costante del moto (costante
    di Carter $Q$) e un sistema ODE a 4 equazioni.

  \item \textbf{Polarizzazione}: i fotoni del disco sono polarizzati per
    scattering dei elettroni (polarizzazione di Thomson). Il piano di polarizzazione
    ruota lungo la geodetica (rotazione di Faraday, effetto gravitazionale di
    Skrotskij).

  \item \textbf{Ombra multipla}: includere i photon ring di ordine superiore
    (immagini fantasma del disco viste dopo 1, 2, 3 orbite attorno alla fotosfera).

  \item \textbf{Plasma nel disco}: la struttura turbolenta del disco può essere
    simulata con MHD (magnetoidrodinamica) procedurale, usando noise di Perlin
    o simplex tridimensionale invece della semplice modulazione sinusoidale.

  \item \textbf{Jet relativistico}: i buchi neri rotanti spesso emettono jet
    di plasma lungo l'asse di rotazione a velocità $\sim 0.99c$. La simulazione
    visiva dei jet richiede l'equazione di trasferimento radiativo per un plasma
    relativisticamente beamizzato.
\end{enumerate}

% ═════════════════════════════════════════════════════════════════════════════
%  APPENDICE
% ═════════════════════════════════════════════════════════════════════════════
\appendix

\section{Tabella dei Simboli}

\begin{table}[h!]
  \centering
  \renewcommand{\arraystretch}{1.4}
  \begin{tabular}{cll}
    \toprule
    \textbf{Simbolo} & \textbf{Nome} & \textbf{Unità (SI)} \\
    \midrule
    $\gmunu$          & Tensore metrico & adimensionale \\
    $\Gmunu$          & Tensore di Einstein & m$^{-2}$ \\
    $\Rmunu$          & Tensore di Ricci & m$^{-2}$ \\
    $\Tmunu$          & Tensore energia-impulso & Pa = J/m$^3$ \\
    $\RS$             & Raggio di Schwarzschild & m \\
    $\Rph$            & Raggio fotosfera & m \\
    $\Risco$          & Raggio ISCO & m \\
    $\bc$             & Parametro d'impatto critico & m \\
    $u = 1/r$         & Inverso del raggio (Binet) & m$^{-1}$ \\
    $\varphi$         & Angolo azimutale & rad \\
    $b = L/E$         & Parametro d'impatto & m \\
    $\mathcal{D}$     & Fattore Doppler relativistico & adimensionale \\
    $\beta = v/c$     & Velocità ridotta & adimensionale \\
    $\gamma$          & Fattore di Lorentz & adimensionale \\
    $T_H$             & Temperatura di Hawking & K \\
    $\dot{M}$         & Tasso di accrescimento & kg/s \\
    $\eta$            & Efficienza di accrescimento & adimensionale \\
    \bottomrule
  \end{tabular}
  \caption{Principali simboli usati nel documento.}
\end{table}

\section{Unità Geometrizzate}

In questo documento usiamo unità geometrizzate in cui:
\begin{equation}
  G = c = 1
\end{equation}
In queste unità ogni grandezza fisica si esprime in potenze di massa $M$:
\begin{align}
  \text{lunghezza} &= GM/c^2 \approx 1477\ \text{m}\ (M/M_\odot) \\
  \text{tempo}     &= GM/c^3 \approx 4.93\ \mu\text{s}\ (M/M_\odot) \\
  \text{velocità}  &= c = 1
\end{align}

Per convertire in SI: moltiplicare per $GM/c^2$ per le lunghezze e $GM/c^3$
per i tempi.

\end{document}
